\documentclass[UTF8,11pt]{ctexart}
\usepackage[a4paper,margin=24mm]{geometry}
\usepackage{amsmath,amssymb,amsthm,mathtools,mathrsfs}
\usepackage[all]{xy}
\usepackage{xurl,hyperref,enumitem}
\usepackage{tikz}
\usetikzlibrary{arrows.meta,cd,decorations.markings,calc,patterns}
\hypersetup{hidelinks,breaklinks=true}
\setlength{\emergencystretch}{3em}
\allowdisplaybreaks
\newtheorem*{theorem}{Theorem}
\newtheorem*{thm}{Theorem}
\newtheorem*{lemma}{Lemma}
\newtheorem*{proposition}{Proposition}
\newtheorem*{prop}{Proposition}
\newtheorem*{corollary}{Corollary}
\newtheorem*{cor}{Corollary}
\newtheorem*{claim}{Claim}
\theoremstyle{definition}
\newtheorem*{definition}{Definition}
\newtheorem*{defn}{Definition}
\newtheorem*{remark}{Remark}
\newtheorem*{example}{Example}
\newcommand{\R}{\mathbb R}
\newcommand{\C}{\mathbb C}
\newcommand{\Z}{\mathbb Z}
\newcommand{\Q}{\mathbb Q}
\newcommand{\N}{\mathbb N}
\newcommand{\F}{\mathbb F}
\newcommand{\D}{\mathbb D}
\newcommand{\bB}{\mathbb B}
\newcommand{\bC}{\mathbb C}
\newcommand{\bR}{\mathbb R}
\newcommand{\bZ}{\mathbb Z}
\newcommand{\bN}{\mathbb N}
\newcommand{\bQ}{\mathbb Q}
\newcommand{\bD}{\mathbb D}
\newcommand{\bF}{\mathbb F}
\newcommand{\CP}{\mathbb{CP}}
\newcommand{\RP}{\mathbb{RP}}
\newcommand{\sO}{\mathscr O}
\newcommand{\sI}{\mathscr I}
\newcommand{\sF}{\mathscr F}
\newcommand{\sG}{\mathscr G}
\newcommand{\sH}{\mathscr H}
\newcommand{\sR}{\mathscr R}
\newcommand{\sM}{\mathscr M}
\newcommand{\sV}{\mathscr V}
\newcommand{\sU}{\mathscr U}
\DeclareMathOperator{\Hom}{Hom}
\DeclareMathOperator{\Ext}{Ext}
\DeclareMathOperator{\Tor}{Tor}
\DeclareMathOperator{\Spec}{Spec}
\DeclareMathOperator{\Proj}{Proj}
\DeclareMathOperator{\supp}{supp}
\DeclareMathOperator{\im}{im}
\DeclareMathOperator{\id}{id}
\DeclareMathOperator{\Tr}{Tr}
\DeclareMathOperator{\tr}{tr}
\DeclareMathOperator{\rank}{rank}
\DeclareMathOperator{\ord}{ord}
\DeclareMathOperator{\dist}{dist}
\DeclareMathOperator{\End}{End}
\DeclareMathOperator{\Aut}{Aut}
\DeclareMathOperator{\Gal}{Gal}
\DeclareMathOperator{\vol}{vol}
\DeclareMathOperator{\Cl}{Cl}
\DeclareMathOperator{\coker}{coker}
\DeclareMathOperator{\codim}{codim}
\DeclareMathOperator{\divergence}{div}
\DeclareMathOperator{\Ric}{Ric}
\DeclareMathOperator{\grad}{grad}
\DeclareMathOperator{\Hess}{Hess}
\newcommand{\abs}[1]{\left\lvert#1\right\rvert}
\newcommand{\sabs}[1]{\lvert#1\rvert}
\newcommand{\babs}[1]{\bigl\lvert#1\bigr\rvert}
\newcommand{\Babs}[1]{\Bigl\lvert#1\Bigr\rvert}
\newcommand{\bbabs}[1]{\biggl\lvert#1\biggr\rvert}
\newcommand{\BBabs}[1]{\Biggl\lvert#1\Biggr\rvert}
\newcommand{\norm}[1]{\left\lVert#1\right\rVert}
\newcommand{\snorm}[1]{\lVert#1\rVert}
\newcommand{\bnorm}[1]{\bigl\lVert#1\bigr\rVert}
\newcommand{\Bnorm}[1]{\Bigl\lVert#1\Bigr\rVert}
\newcommand{\bbnorm}[1]{\biggl\lVert#1\biggr\rVert}
\newcommand{\BBnorm}[1]{\Biggl\lVert#1\Biggr\rVert}
\newcommand{\widebar}[1]{\overline{#1}}
\newcommand{\nicefrac}[2]{\frac{#1}{#2}}
\newcommand{\linnprod}[2]{\langle#1,#2\rangle}
\newcommand{\myindex}[1]{#1}
\newcommand{\glsadd}[1]{}
\newcommand{\avoidbreak}{}
\newcommand{\sourceref}[1]{\textnormal{[原文：\nolinkurl{#1}]}}
\newcommand{\thmref}[1]{\sourceref{#1}}
\newcommand{\lemmaref}[1]{\sourceref{#1}}
\newcommand{\propref}[1]{\sourceref{#1}}
\newcommand{\corref}[1]{\sourceref{#1}}
\newcommand{\defnref}[1]{\sourceref{#1}}
\newcommand{\exampleref}[1]{\sourceref{#1}}
\newcommand{\exerciseref}[1]{\sourceref{#1}}
\newcommand{\figureref}[1]{\sourceref{#1}}
\newcommand{\sectionref}[1]{\sourceref{#1}}
\newcommand{\appendixref}[1]{\sourceref{#1}}
\newcommand{\stacksref}[2]{\href{https://stacks.math.columbia.edu/tag/#1}{#2 [#1]}}


\newcommand{\E}{\mathbb E}
\newcommand{\Prob}{\mathbb P}
\newcommand{\Reff}{\mathcal R_{\mathrm{eff}}}
\newcommand{\eps}{\varepsilon}
\DeclareMathOperator{\diam}{diam}
\DeclareMathOperator{\Var}{Var}
\DeclareMathOperator{\Crit}{Crit}
\newcommand{\dd}{\,\mathrm d}


\begin{document}
\section*{085\quad 纤维丛上同调的Leray--Hirsch模分解}
\subsection*{来源与整理范围}
Haynes Miller 的 MIT 18.906 课程，Sanath Devalapurkar 等现场 TeX 记录，2016--2017；Lecture 68，Leray--Hirsch 小节及其紧接在前的滤过说明。仓库提交 \nolinkurl{3f5d3189e2082716a69fccc1711d02ed848552d2}。
\par\url{https://github.com/sanathdevalapurkar/algtop-notes/blob/3f5d3189e2082716a69fccc1711d02ed848552d2/906/lec-68-leray-hirsch.tex}
\par\textbf{恢复方式（整理者说明）：}依据人类原文重新整理以补齐缺失题号；并非旧题证文件的逐字节恢复；2026-09-28。
\subsection*{作者命题与人类证明}

\paragraph{Filtration used in the proof.}
Let $\pi:E\to B$ be a fibration, with path-connected base and fibre $F$.
Cohomology has coefficients in a commutative ring $R$. The multiplicative Serre
spectral sequence is
\[
 E_2^{s,t}=H^s(B;\underline{H^t(F;R)})\Longrightarrow H^{s+t}(E;R).
\]
The pullback $\pi^*$ makes $H^*(E;R)$ an $H^*(B;R)$-module. On each total
degree, rewrite the decreasing Serre filtration as
$F_tH^n(E;R)=F^{n-t}H^n(E;R)$. Then
\[
 F_0H^n(E;R)=E_\infty^{n,0}=\im(\pi^*:H^n(B;R)\to H^n(E;R)),
 \qquad \operatorname{gr}_tH^*(E;R)=E_\infty^{*,t}.
\]
This is an exhaustive, bounded-below filtration by $H^*(B;R)$-modules, and
the displayed identification is an identification of modules.

\begin{theorem}[Leray--Hirsch]
Suppose $B$ and $F$ are path connected, each $H^t(F;R)$ is a finite-rank free
$R$-module, and restriction $i^*:H^*(E;R)\to H^*(F;R)$ is surjective.
Choose a degree-preserving $R$-linear section
$\sigma:H^*(F;R)\to H^*(E;R)$ of $i^*$. Then
\[
 \overline\sigma:H^*(B;R)\otimes_RH^*(F;R)\longrightarrow H^*(E;R),
 \qquad x\otimes y\longmapsto\pi^*(x)\smile\sigma(y)
\]
is an isomorphism of $H^*(B;R)$-modules.
\end{theorem}
\begin{remark}
The isomorphism depends on the choice of $\sigma$. In particular, the theorem
asserts freeness as an $H^*(B;R)$-module, not a canonical algebra isomorphism.
\end{remark}
\begin{proof}
Use the Serre spectral sequence above. Restriction to the fibre is an edge
homomorphism and factors through
\[
 H^*(E;R)\longrightarrow E_\infty^{0,*}
 \hookrightarrow E_2^{0,*}=H^0(B;\underline{H^*(F;R)})
 \hookrightarrow H^*(F;R).
\]
Surjectivity of restriction implies
$H^0(B;\underline{H^*(F;R)})=H^*(F;R)$. Thus the monodromy action of
$\pi_1(B)$ on $H^*(F;R)$ is trivial. Since each graded piece of fibre
cohomology is finite free, the second page is
\[
 E_2^{s,t}=H^s(B;R)\otimes_R H^t(F;R).
\]
There are no outgoing differentials on the fibre column: otherwise the
restriction to the fibre could not be surjective. There are no outgoing
differentials on the base row either. The differential is a derivation,
and this tensor-product algebra is generated by the base row and fibre
column. Consequently all differentials vanish, page by page. Hence
\[
 E_
\infty^{*,t}=H^*(B;R)\otimes_RH^t(F;R).
\]
Filter the tensor product by degree in fibre cohomology:
\[
 F_q\bigl(H^*(B;R)\otimes_RH^*(F;R)\bigr)
   =\bigoplus_{t\leq q}H^*(B;R)\otimes_RH^t(F;R).
\]
The map $\overline\sigma$ preserves this filtration and the corresponding
filtration on $H^*(E;R)$. On the associated graded its map is the
identification
\[
 H^*(B;R)\otimes_RH^t(F;R)
   =E_\infty^{*,t}=\operatorname{gr}_tH^*(E;R).
\]
It is therefore an isomorphism on the associated graded. Since the
filtrations are exhaustive and bounded below, $\overline\sigma$ itself is
an isomorphism.
\end{proof}

\subsection*{整理说明与先修范围（非作者正文）}
按原证明保留局部系数、边同态、乘法微分、滤过与伴随分次的完整论证。删除与定理无关的课程通知及 Dress 谱序列旁述；统一宏和口语句式，属编辑转录。

原定理把 $\sigma$ 误称为 surjection，此处按其作为限制映射右逆的用途改为 section，并明确保持次数；原文“$\pi_1(B)$ 在 $F$ 上作用平凡”按对应公式改为在 $H^*(F;R)$ 上作用平凡。结论限于模同构，不提升为环同构。

先修是带乘法的 Serre 谱序列、边同态和有限自由系数的上同调换系数。难点在于从全局类的可延拓性控制单值作用及全部微分，再用所选提升把伴随分次分解实现为总上同调的实际模同构。
\subsection*{可修改的简短标注（非作者正文）}
$c$：纤维丛总空间的上同调结构

$a$：Serre谱序列；局部系数与单值作用；杯积；边同态；滤过与伴随分次

\texttt{cross\_theory\_status = yes}。将纤维上类能提升的几何信息转化为谱序列中微分消失，再用杯积构造返回总空间的模同构；见边同态分解及最后的滤过比较。

全部标注为非穷尽、可修改初稿。
\subsection*{来源许可与编辑归属}
上游仓库 LICENSE 为 MIT License，Copyright (c) 2017 Sanath。作者与许可全文见本包 \nolinkurl{sources/licenses/MIT_algtop_notes.txt}。
ChatGPT 加入题号、中文来源说明、整理范围和初稿标注；编辑说明与人类证明分列。
\end{document}
